\section*{Chapter \ref{ch:g11:monopoly} --- Monopoly}
% figdata: figdata/grade-11/04-monopoly/out/generics.csv (cited, ledger g11:fda2019-*); Brenland numbers invented.

\begin{solution}{exo:g11:monopoly:1}
$R = 50Q - 2Q^2$, $R' = 50 - 4Q$; at $Q = 5$, $R' = 30$ (the price there is 40).
\end{solution}

\begin{solution}{exo:g11:monopoly:2}
Revenue 10, 18, 24, 28, 30; \omterm{def:g11:marginal-reasoning:marginal-benefit}{marginal revenue} 10, 8, 6, 4, 2.
\end{solution}

\begin{solution}{exo:g11:monopoly:3}
$60 - 2Q = 10$ gives $Q = 25$, $p = 35$, profit $(35 - 10) \times 25 = 625$.
\end{solution}

\begin{solution}{exo:g11:monopoly:4}
61.4\,\%, 46.5\,\%, 21.2\,\% and 6.1\,\% of the brand's price. The first producer's price is 38.6\,\% below the brand's; the second brings it down by $(61.4 - 46.5)/61.4 \approx
24.3\,\%$; the fourth by $(32.2 - 21.2)/32.2 \approx 34.2\,\%$; the sixth by $(14.4 - 6.1)/14.4 \approx 57.6\,\%$.
\end{solution}

\begin{solution}{exo:g11:monopoly:5}
\omterm{def:g11:monopoly:markup}{Markup} $50/30 \approx 1.67$; \omterm{def:g11:monopoly:markup}{Lerner index} $20/50 = 0.4$. By \cref{prop:g11:monopoly:lerner}, $|\varepsilon| = 1/0.4 = 2.5$: $\varepsilon = -2.5$.
\end{solution}

\begin{solution}{exo:g11:monopoly:6}
The competitive output is 50 (price 10). \omterm{def:g10:what-prices-do:dwl}{Deadweight loss} $\tfrac12 \times (50 - 25) \times (35 - 10) = 312.5$.
\end{solution}

\begin{solution}{exo:g11:monopoly:7}
$100 - 4Q = 20$ gives $Q = 20$, $p = 60$. Competitive: $100 - 2Q = 20$, $Q = 40$. \omterm{def:g10:what-prices-do:dwl}{Deadweight loss} $\tfrac12 \times 20 \times 40 = 400$.
\end{solution}

\begin{solution}{exo:g11:monopoly:8}
With $\mathit{MC} = 30$: $100 - 2Q = 30$, $Q = 35$, $p = 65$, a rise of 5 for a cost rise of 10. In general $a - 2bQ = c$ gives $Q = (a - c)/(2b)$ and $p = a - (a - c)/2 = (a + c)/2$:
the price rises by half of any rise in $c$.
\end{solution}

\begin{solution}{exo:g11:monopoly:9}
Different prices: in A, $40 - 2Q_A = 8$, $Q_A = 16$, $p_A = 24$; in B, $24 - 2Q_B = 8$, $Q_B = 8$, $p_B = 16$; profit $16 \times 16 + 8 \times 8 = 320$. One price: below 24 both
markets buy, $Q = 64 - 2p$, so $p = 32 - Q/2$ and $\mathit{MR} = 32 - Q = 8$ gives $Q = 24$, $p = 20$ (A buys 20, B buys 4); profit $12 \times 24 = 288$ (selling in A only, at 24,
gives 256). Discrimination raises the profit by 32 with the same total output.
\end{solution}

\begin{solution}{exo:g11:monopoly:10}
At \omterm{def:g11:marginal-reasoning:marginal-benefit}{marginal cost}, $p = 2$ and $Q = 2000$: revenue 4000, cost 7600, a loss of 3600, the fixed cost. It just covers its costs where $p = 2 + 3600/Q$ and $p = 22 - Q/100$: at $Q =
1800$, $p = 4$ (check: $2 + 3600/1800 = 4$ and $22 - 18 = 4$).
\end{solution}

\begin{solution}{exo:g11:monopoly:11}
Manufacturer prices: $(0.614 - 0.465)/0.614 \approx 24.3\,\%$. Pharmacy prices: $(0.696 - 0.562)/0.696 \approx 19.3\,\%$. Pharmacies buy through wholesalers, whose \omterm{def:g11:monopoly:markup}{markup} is
added to the manufacturer's price; it makes invoice prices higher and falls less in proportion when the manufacturers' prices fall.
\end{solution}

\begin{solution}{exo:g11:monopoly:12}
A price at which demand is inelastic, since a monopoly never sells there; a \omterm{def:g11:monopoly:markup}{markup} far from one over the elasticity; or prices that do not fall, or rise, after the patent
expires and generic makers enter. Before concluding, one must check that the firm is free to set its price (it may be regulated or negotiated with insurers) and that the
elasticity is measured at the right price.
\end{solution}

\begin{solution}{exo:g11:monopoly:13}
Fares close to the average cost of flying the route, since a higher fare would bring an entrant that could undercut it and leave at no loss. The prediction weakens when
entry carries sunk costs (airport slots, advertising, ground staff), when the incumbent can cut fares as soon as an entrant appears and raise them once it leaves, or when
slots at the airports are scarce.
\end{solution}

\begin{solution}{exo:g11:monopoly:14}
By \cref{prop:g11:monopoly:lerner}, $\mathit{MR} = p(1 - 1/|\varepsilon|)$, which is negative when $|\varepsilon| < 1$. At such a quantity \omterm{def:g11:marginal-reasoning:marginal-benefit}{marginal revenue} cannot equal a positive
\omterm{def:g11:marginal-reasoning:marginal-benefit}{marginal cost}, and selling one unit fewer raises revenue and saves its cost: profit rises. So the monopoly's quantity is never one where demand is inelastic.
\end{solution}

\begin{solution}{exo:g11:monopoly:15}
Before: $Q = 80$, $p = 20$, \omterm{def:g10:surplus-and-efficiency:consumer-surplus}{consumer surplus} $\tfrac12 \times 80 \times 80 = 3200$, \omterm{def:g10:surplus-and-efficiency:producer-surplus}{producer surplus} 0, total 3200. After: $100 - 2Q = 10$, $Q = 45$, $p = 55$; \omterm{def:g10:surplus-and-efficiency:consumer-surplus}{consumer surplus} $\tfrac12
\times 45 \times 45 = 1012.5$; \omterm{def:g10:surplus-and-efficiency:producer-surplus}{producer surplus} $(55 - 10) \times 45 = 2025$; total 3037.5. \omterm{def:g10:surplus-and-efficiency:total-surplus}{Total surplus} falls by 162.5 and consumers lose 2187.5. The cost saving does not make up
for the lost competition: the merger fails both tests, since \omterm{def:g10:surplus-and-efficiency:total-surplus}{total surplus} falls by 162.5 and \omterm{def:g10:surplus-and-efficiency:consumer-surplus}{consumer surplus} by 2187.5.
\end{solution}

\begin{solution}{pb:g11:monopoly:1}
Fares in dalens; crossings a day.

\textbf{1.} $p = 24 - Q/100$.

\textbf{2.} $R = 24Q - Q^2/100$; $R' = 24 - Q/50$.

\textbf{3.} $24 - Q/50 = 4$ gives $Q = 1000$; $p = 14$.

\textbf{4.} $(14 - 4) \times 1000 - 3600 = 6400$ dalens a day.

\textbf{5.} $\varepsilon = -100 \times 14/1000 = -1.4$; $L = (14 - 4)/14 \approx 0.714 = 1/1.4$.

\textbf{6.} $24 - Q/100 = 4$ gives $Q = 2000$.

\textbf{7.} $\tfrac12 \times (2000 - 1000) \times (14 - 4) = 5000$ dalens a day.

\textbf{8.} Monopoly: $\tfrac12 \times 1000 \times (24 - 14) = 5000$. At 4: $\tfrac12 \times 2000 \times 20 = 20000$.

\textbf{9.} Of the 20000 (before the fixed cost), travellers get 5000, the company 10000, and 5000 are lost.

\textbf{10.} At 4 it covers only the cost of each crossing and loses its fixed cost of 3600 a day; at 14 it makes 6400.

\textbf{11.} $0 - 3600 = -3600$ dalens a day.

\textbf{12.} At $p = 6$, $Q = 2400 - 600 = 1800$, and average cost is $4 + 3600/1800 = 6$: the company breaks even.

\textbf{13.} $\tfrac12 \times (2000 - 1800) \times (6 - 4) = 200$.

\textbf{14.} $(5000 - 200)/5000 = 96\,\%$.

\textbf{15.} The passes bring in $1200 \times 3 = 3600$ a day, the fixed cost, and each crossing pays its cost: profit 0. Crossings are 2000 a day and the \omterm{def:g10:what-prices-do:dwl}{deadweight loss} is 0.

\textbf{16.} The ferry's profit rises with its fare up to 14, but above 8 it loses every traveller to the bridge. It charges 8, makes 1600 crossings, and earns $4 \times 1600 - 3600 =
2800$ a day.

\textbf{17.} $\tfrac12 \times (2000 - 1600) \times (8 - 4) = 800$.

\textbf{18.} The ferry may keep its fare at or below 8 so that the bridge is never built: the threat of entry alone limits its price.

\textbf{19.} The merged firm is a monopoly again; with the ferry's costs it charges 14. With a \omterm{def:g11:marginal-reasoning:marginal-benefit}{marginal cost} $c \geq 0$ its price would be $(24 + c)/2 \geq 12$, above 8, so no cost
saving on crossings could make the merger good for travellers here.

\textbf{20.} The ferry monopoly's \omterm{def:g10:what-prices-do:dwl}{deadweight loss} is 5000 dalens a day; a fare regulated at average cost, 6 dalens, cuts it to 200 and so recovers 96\,\% of it.
\end{solution}
